\begin{center}
\begin{tabular}{lrrrr}
\toprule
Integer quartic & Variables & Rows & Terms & Height\\
\midrule
Matched incoming parent & 517 & 310 & 6,774 & 3,904,267\\
Cofinal/modulus checkpoint & 512 & 310 & 6,742 & 3,904,268\\
\textbf{Signed-factorial principal} & \textbf{508} & \textbf{308}
 & \textbf{6,545} & \textbf{996,521}\\
\addlinespace
Incoming lower-height alternative & 519 & 321 & 6,253 & 411,060\\
\bottomrule
\end{tabular}
\end{center}
All three have exact degree four. The 512-variable checkpoint slightly
increases height; the principal 508-variable equation improves every
displayed quantity relative to the matched parent. A separate family has
ten strictly positive unknowns, with operation/degree pairs
$89/288$, $91/283$, $92/278$, $94/262$, $96/258$, and $97/254$.
Section~\ref{sec:positive} specifies those polynomials and their domains.
The separate 519-variable incoming equation retains fewer terms and smaller
coefficients than the 508-variable result; it remains a useful tradeoff.
